\subsection{A family with zero compact abelianization}

The preceding construction detects noncuspidality through a nontorsion
compact class. A different family shows that the obstruction can be
entirely in the boundary part of geometric homology.

\begin{lemma}[Lifting over a source of $p$-cohomological dimension one]
\label{lem:cdlift}
Let $\Gamma_0$ be profinite with $\operatorname{cd}_p\Gamma_0\leq1$,
and let $N_0$ be a finitely generated pro-$p$ group with a continuous
$\Gamma_0$-action. Every continuous cocycle
$\Gamma_0\to N_0^{\ab}$ lifts to a continuous $N_0$-valued cocycle.
\end{lemma}
\begin{proof}
Pull $N_0\rtimes\Gamma_0\to N_0^{\ab}\rtimes\Gamma_0$ back along
the section specified by the cocycle. The resulting extension
$E\to\Gamma_0$ has kernel $D=[N_0,N_0]$.
Choose a nested cofinal sequence $V_i$ of characteristic open subgroups
of $N_0$. Then $E_i=E/(D\cap V_i)$ has finite $p$-group kernel and
$E=\varprojlim_iE_i$.

Every extension of $\Gamma_0$ by a finite $p$-group $L$ splits.
To see this, induct on $|L|$. The nontrivial characteristic elementary
abelian subgroup $Z=\Omega_1(Z(L))$ is normal in the extension.
First split the quotient extension with kernel $L/Z$ by induction,
then pull back along its section. The remaining extension has kernel
$Z$, and its obstruction lies in $H^2(\Gamma_0,Z)=0$.
The action on $Z$ is the one induced by this extension. Continuous
cohomology applies: an extension with finite kernel has local
continuous sections, by choosing an open subgroup meeting its kernel
trivially, and hence a continuous set-theoretic section on finitely
many clopen translates.

If a section into $E_i$ has been chosen, pull $E_{i+1}\to E_i$ back
along it. The same finite-$p$-kernel argument supplies a compatible
section into $E_{i+1}$. Taking this sequence's inverse limit gives a
continuous section $\Gamma_0\to E$. Its $N_0$-component lifts the
specified abelian cocycle exactly.
\end{proof}

\begin{proposition}\label{prop:boundaryfamily}
For every $n=3^r$, $r\geq1$, there are continuum many pairwise
geometrically nonconjugate full pro-$3$ sections of $U_n/\Q$ whose
whole cocycles are unramified outside $3$ and whose classes in
$H^1(\Q,T_3\Jac(X_n))$ are zero. Each fails point-or-cusp geometry
at $3$. All these assertions persist after passing to any
Galois-stable geometric quotient retaining $N_n^{\ab}$.
\end{proposition}
\begin{proof}
Use $x=X/Z$, $y=Y/Z$ and the standard tangent $t=x$ at
$Q_0=[0:-1:1]$. In its Puiseux fiber functor,
\[
x=t,\qquad y=-(1+t^n)^{1/n}.
\]
The leading constants of $x,y$ are $1,-1$; both have compatible
rational $3$-power roots. The reference section therefore has zero
Kummer classes under these two invertible functions.

Set $M=N_n^{\ab}$. The rational cusps $Q_0$ and
$Q_1=[-1:0:1]$ supply Galois-equivariant maps
$j_0,j_1:\Z_3(1)\to M$ from their inertia groups.
The map induced by the functions $x,y$ satisfies
\[
(q_x,q_y)j_0(z)=(z,0),\qquad
(q_x,q_y)j_1(z)=(0,z),
\]
since these are their orders of zero at the two cusps.
All boundary classes have zero image in $T_3\Jac(X_n)$.

Let $\Gamma=\Gamma_{\{3\}}$ be the arithmetic source above.
The free subgroup $P$ has index two, so
$\operatorname{cd}_3\Gamma=1$: the Hochschild--Serre spectral sequence
and invertibility of two on every $3$-primary module reduce the
assertion to $\operatorname{cd}_3P=1$.
The Kummer cocycle $\kappa_3(3)$ factors through $\Gamma$.
Indeed the fields obtained by adjoining compatible $3$-power roots
of $3$ and roots of unity are unramified outside $3$, and over
$K=\Q(\zeta_3)$ their Galois closures are pro-$3$.
For every $b\in3+9\Z_3$ the abelian cocycle
\[
a_b=(j_0+bj_1)\kappa_3(3)
\]
therefore lives on $\Gamma$. Lemma~\ref{lem:cdlift} lifts it to an
actual $N_n$-cocycle on $\Gamma$. Inflate to $G_\Q$.
The entire cocycle is unramified outside $3$ because this is true
of $G_\Q\to\Gamma$, and its compact abelianization is zero.

Localize at $3$ and apply the valuation to the two coordinate Kummer
classes. The resulting vector is $(1,b)$, so different parameters
give different geometric conjugacy classes.
If a point $(x,y)\in U_n(\Q_3)$ gave this vector, then
$v_3(x)=1$. The equation $y^n=-1-x^n$ would imply $v_3(y)=0$,
contradicting $b\in3+9\Z_3$.
The only rational cusps over $\Q_3$ are $Q_0,Q_1$ and
$Q_\infty=[1:-1:0]$, since $\Q_3$ has no nontrivial
$3$-power roots of unity. Their local coordinate orders give,
respectively, the three lines of possible cusp valuation vectors
\[
(z,0),\qquad(0,z),\qquad(-z,-z)\qquad(z\in\Z_3).
\]
The vector $(1,b)$ is on none of them: $b$ is nonzero and is not one.
Thus the sections are neither point geometric nor cusp geometric
at $3$. This also proves global noncuspidality. The coordinate maps
and all these valuation distinctions factor through $M$, so they
remain valid in every stated quotient.
\end{proof}

The two constructions exhibit different failures of local geometry.
The first is visible in the proper elliptic Tate module; the second
vanishes in compact homology and is detected by two boundary Kummer
coordinates. Neither construction gives a Selmer counterexample.
